
\chapter*{Prefácio}
\addcontentsline{toc}{chapter}{Prefácio}
\markboth{Prefácio}{Prefácio}

\section*{A quem se dirige}

O relatório serve a quem precisa de uma estimativa de exposição ao corte antes
de decidir sobre um ativo fotovoltaico, e a quem modela geração renovável a
partir de irradiância ou de imagem de satélite. Para o primeiro leitor, os
capítulos de despacho e de rede dizem quanto se corta e o que prediz onde se
corta. Para o segundo, o capítulo~\ref{cap:ancora} estabelece que a geração
verificada carrega uma decisão operativa, e não apenas o recurso físico.

Nenhum capítulo pressupõe acesso a dado sob contrato. Todas as fontes são
abertas e sem autenticação, e o apêndice~\ref{ap:procedencia} nomeia cada uma
com o pacote de origem.

\section*{O que fica de fora}

\campo{Micro e minigeração distribuída} O regime de restrição da geração
distribuída é distinto do da geração centralizada, e o corte apurado pelo
operador não a cobre. O capítulo~\ref{cap:satelite} mede a detectabilidade
dessas instalações por satélite, e não o corte delas.

\campo{Ressarcimento} A Lei 15.269/2025 criou compensação por
\textit{constrained-off} durante a execução das análises. O relatório valora a
energia cortada ao Custo Marginal de Operação, que mede valor sistêmico e não
ressarcimento. Os dois divergem por construção.

\campo{Decisão de investimento} O capítulo~\ref{cap:aplicacao} estima quanto
uma área nova produziria e quanto o despacho absorveria. Contrato, outorga e
parecer de acesso não estão em dado aberto e não entram na estimativa.

\campo{Corte anterior a abril de 2024} A apuração pública de
\textit{constrained-off} fotovoltaico começa nessa data. O período anterior é
tratado por fator de capacidade, que responde a várias causas ao mesmo tempo.
